

\section{The Dirichlet process}
\label{sec:ch4_Dirichilet_process}

\relatedrefs{
    The presentation in this chapter (to be removed) follows \cite[Ch.~1-2]{Orbanz2014BayesianNP}
}



\paragraph{Bayesian mixture model} We are very familiar at this stage with the hierarchical GMM, let us consider a different view of this model. Denoting the assignment variable as $z\in\N$, the distribution the observation $x$ generated by a the $k$-th cluster can be expressed as 
\begin{equation}
    p_k(x) = p(x \mid z=k),
\end{equation}
where we have not assumed that this distribution is Gaussian. Furthermore, the probability for a given observation to be generated by such cluster, can be denotes as 
\begin{equation}
    g_k = \Prob{Z=k},
\end{equation}
where $\sum_{\in\N}c_k = 1, c_k\geq0,\, \forall k\in\N$. The resulting model, given by
\begin{equation}
p(x) = \sum_{k\in\N} c_k p(x \mid z=k).    
\end{equation}
We will refer to this distribution as a mixture model. Furthermore, when there is only a finite number of probabilities $c_k$, we will say that this is a finite mixture. 

The set of all sequences of the form $\{c_k\}_{k\in\N}$ is called \emph{simplex}, denoted by 
\begin{equation}
    \Delta \defeq \left\{ \{c_k\}_{k\in\N} \middle| \sum_{k\in\N}c_k = 1, c_k\geq0 \right\}.
\end{equation}
We can also assume that the cluster densities are in a parametric model over $\cX$, given by $\{p(\cdot \mid \phi) | \phi\in\Phi\}$. This way, the assignment variable informs the chosen parameters for the conditional probabilites. This yields the following expression for the generative model
\begin{equation}
    p(x) = \sum_{k\in\N} c_k p(x \mid \phi).    
\end{equation}
This allows us to express the following interpretation: Consider a discrete (atomic) probability measure $\theta$ on the parameter space $\Phi$, given by:
\begin{equation}
    \theta(\cdot) = \sum_{k\in\N} c_k \delta_{\phi_k}(\cdot). 
\end{equation}
This measure assigns probability to atoms in the parameter space and is referred to as the mixing measure---see Fig.~\ref{fig:mixing_model}. This is because we can express the PGM as
\begin{align}
    p(x) 
    &= \int_\Phi p(x \mid \phi) \theta(\phi)\d\phi\\
    &= \int_\Phi p(x \mid \phi) \sum_{k\in\N} c_k \delta_{\phi_k}(\phi)\d\phi\\
    &= \sum_{k\in\N}  c_k  p(x \mid \phi_k).
\end{align}

\begin{remark}
    The statistical model used in clustering can be written as the mixture model above. This implies that the parameter of such model is a discrete probability distribution (called the mixing distribution). 
\end{remark}


\begin{figure}
    \centering
    \includegraphics[width=0.7\textwidth]{img/week4_mixing_measure.pdf}
    \caption{Illustration of a mixing measure (left) and the mixing model (right).}
    \label{fig:mixing_model}
\end{figure}



Having identified the parameter space of the mixture model, let us now consider a Bayesian mixture model, that is, a random mixing measure 
\begin{equation}
    \Theta = \sum_{k\in\N} C_k \delta_{\Phi_k}.
\end{equation}
Therefore, the prior of a Bayesian mixture is the distribution of the random mixing law $\Theta$. Constructing this prior is not difficult. First, we can to define a prior for the parameters of the components, that is, 
\begin{equation}
    \Phi_1,\Phi_2,\ldots, \sim_\text{i.i.d.} G,
\end{equation}
and independent of $C_k$. Then to sample the weights we cannot consider indepenence, that is the collection elements of the sequence $\{c_k\}_{k\in\N}$ needs to add up to one. However, in the finite case, we can simply sample $K$ i.i.d.~elements $V_k$ form [0,1], and then define
\begin{equation}
    C_k \defeq \frac{V_x}{V_1+\ldots+V_K}.    
\end{equation}

\begin{remark}
    When the RVs $V_k$ follow a gamma distribution, $C_k$ follow a Dirichlet distribution. 
\end{remark}

\paragraph{Stick-breaking construction.} The above construction works well for finite $k$ but very often we will want to define the mixture over infinite components (Why?). For an infinite number of components, the sum of i.i.d.~variables $V_1 + V_2 + \ldots$ will diverge almost surely. 

Sampling an infinite number of parameters $\Phi_k\sim G$ is not problematic. For the mixing weights $C_k$, we can consider the stick-breaking construction, which provides a simple way to generate an infinite sequence of non-negative weights that sum to one.

Imagine a stick of length 1 representing total probability mass. The idea is to break this stick an infinite number of times, where the remaining pieces of the stick will be the $c_k$.

At step $k$, we:
\begin{enumerate}
  \item break off a fraction $V_k \in (0,1)$ of the remaining stick,
  \item assign this piece length $C_k$ to component $k$,
  \item keep the rest for future components.
\end{enumerate}

Formally, let $|I_1| = 1$. For $k = 1,2,\ldots$:
\[
V_k \sim H, \qquad
C_k = |I_k|\,V_k, \qquad
|I_{k+1}| = (1 - V_k)\,|I_k|.
\]

The weights $(C_k)_{k\ge1}$ are non-negative and satisfy
\[
\sum_{k=1}^\infty C_k = 1 \quad \text{almost surely}.
\]

This construction allows us to define mixture models with infinitely many components, while ensuring that the total probability mass remains finite and well defined.

\begin{figure}
    \centering
    \includegraphics[width=0.9\textwidth]{img/week4_stick_breaking.pdf}
    \caption{Illustration of a the stick breaking process.}
    \label{fig:stick_breaking}
\end{figure}

The stick-breaking procedure allows for constructing a mixing measure with infinite components by choosing a specific distribution for the components's parameters $G$. Choosing $H$ above as a beta distribution we have: 

\begin{definition}[Dirichlet process]
    Consider $\alpha>0$ a concentration parameter,  and $G$ a probability measure on $\Phi$. The random discrete probability measure $\Theta$ generated by 
    \begin{align}
        V_1,V_2,\ldots, &\sim_\text{i.i.d.} \text{Beta}(1,\alpha) \quad \text{and}\quad C_k \defeq V_k\prod_{j=1}^{k-1}(1-V_k)\\
        \Phi_1,\Phi_2,\ldots, &\sim_\text{i.i.d.} G
    \end{align}
    is called a Dirichlet process (DP) with base measure $G$ and concentration $\alpha$.
    
\end{definition}

Fig.~\ref{fig:2DPs} shows draws from two DPs with different concentration parameters. 

\begin{figure}
    \centering
    \includegraphics[width=0.9\textwidth]{img/week4_2DPs.pdf}
    \caption{Draws from two DPs with a standard Gaussian as base measure, and different concentration parameters.}
    \label{fig:2DPs}
\end{figure}





\paragraph{Dirichlet process mixture model.}
Using a Dirichlet process as a prior on the mixing measure yields an infinite mixture model.
Specifically, let
\[
\Theta \sim \mathrm{DP}(\alpha, G), \qquad
\phi_i \mid \Theta \sim \Theta, \qquad
x_i \mid \phi_i \sim p(x \mid \phi_i),
\]
for $i=1,\ldots,N$.
Marginally, this defines a mixture model with a countably infinite number of components, where the number of components effectively used by the data grows with $N$.

\paragraph{Posterior of the Dirichlet process.}
A key property of the Dirichlet process is conjugacy.
If
\[
\Theta \sim \mathrm{DP}(\alpha, G)
\quad\text{and}\quad
\phi_1,\ldots,\phi_N \mid \Theta \sim \Theta,
\]
then the posterior distribution of $\Theta$ is again a Dirichlet process:
\[
\Theta \mid \phi_1,\ldots,\phi_N
\;\sim\;
\mathrm{DP}\!\left(
\alpha + N,\;
\frac{\alpha G + \sum_{n=1}^N \delta_{\phi_n}}{\alpha + N}
\right).
\]

This result shows that Bayesian updating under the Dirichlet process corresponds to combining the prior base measure with the empirical distribution of the observed atoms.

